\section{Material complementario}

Ademas del anexo principal, el trabajo incorpora material complementario destinado a ampliar la documentacion del proceso de investigacion.

Este material puede incluir:

\begin{itemize}
    \item bases de datos brutas o semiprocesadas;
    \item resultados intermedios de depuracion y transformacion;
    \item cartografia adicional;
    \item figuras auxiliares y salidas no incorporadas al cuerpo principal del TFM.
\end{itemize}

La separacion entre anexo principal y material complementario responde a un criterio de claridad expositiva y de gestion tecnica de archivos.
